% ---------------------------------------------------------------------------
%  01-problem.tex -- what the tool is asked to do, and why it is hard
% ---------------------------------------------------------------------------
\section{The problem, and why parallelising it is not routine}
\label{sec:problem}

\subsection{The question the tool answers}

The input is a network of communicating processes with real-valued clocks --
the standard \uppaal{} model -- and a question such as ``can two philosophers
ever eat at the same time?'' or ``can the system ever deadlock?''.  The tool
must answer \emph{yes} or \emph{no}, and the answer must be right.

\begin{intuition}
Think of a board game where several players move at once and each carries a
stopwatch.  A \emph{state} is where all the pieces are plus what all the
stopwatches read.  Because stopwatches are real numbers, there are infinitely
many states, so you cannot just enumerate them.  The trick that makes timed
model checking possible is to lump together all the stopwatch readings that
behave identically into one \emph{zone}, and then explore a finite graph of
(board position, zone) pairs.
\end{intuition}

Formally, the state space explored is the \emph{zone graph}: nodes are pairs
$(\dk, Z)$ where $\dk$ is the discrete part (one location per process plus the
integer variables) and $Z$ is a convex set of clock valuations, and edges are
the symbolic transitions ``take an enabled edge, reset some clocks, let time
pass, re-apply the invariants''.  The search is a fixpoint computation: keep
generating successors until nothing new appears.

\subsection{Why this is not an embarrassingly parallel problem}

Reachability looks parallel: many states, expand them independently.  The
difficulty is the \emph{passed set}, the structure that remembers what has
already been explored.  Three properties of it fight parallelism, and they
compound.

\paragraph{(1) The passed set is not a set of points; it is a set of sets.}
An untimed model checker stores visited states in a hash set: membership is an
exact match, entries are immutable, and a lock-free hash table (\ltsmin{}'s
design) handles it beautifully.  Here, a stored entry is a \emph{zone}, i.e.\
an infinite set of valuations, and the test is not equality but
\emph{coverage}: a new zone can be discarded if some stored zone already
contains it, or -- coarser and much more effective -- \emph{simulates} it.
Coverage is not a hash lookup.

\paragraph{(2) The passed set is mutable.}
When a new zone arrives that covers stored ones, those stored ones are
\emph{deleted}.  When two stored zones happen to have a convex union, they are
\emph{replaced} by a single zone.  So the structure a reader is consulting is
being rewritten underneath it.  A lock-free read of an immutable hash table is
easy; a lock-free read of a structure whose entries are being deleted and
replaced is normally unsafe.

\paragraph{(3) The interesting work is inside the critical section.}
In an untimed checker the critical section is a compare-and-swap.  Here it is
subsumption tests and convex-union tests over $\dim \times \dim$ matrices --
tens of microseconds if done naively.  Putting that behind one lock serialises
the whole tool; putting it behind many locks only helps if the per-lock work is
genuinely small.

\begin{intuition}
Parallel search over a plain visited-set is a solved problem: everyone shares
one big hash table and atomically claims entries.  Here the shared table
contains \emph{regions}, not points, and inserting a region can erase or absorb
regions that are already there.  The whole design problem is: how do you let
$32$ threads read and rewrite a structure like that, concurrently, without
either locking the world or getting a wrong answer?
\end{intuition}

\subsection{The shape of the answer}

The engine's answer has three parts, and each addresses one of the three
difficulties above.

\begin{enumerate}[leftmargin=1.5em, itemsep=3pt]
\item \textbf{Split the shared structure by discrete state.}  The passed set is
      one flat table divided into $\NB \ge 16\,384$ independently locked
      stripes.  Two workers conflict only when the discrete states they are
      inserting hash to the same stripe -- measured at $2$--$3\%$
      (\cref{sec:parallel}).

\item \textbf{Make the common case a \emph{shared} read.}  $83\%$ of insertion
      attempts turn out to be ``already covered, drop it''.  That test is
      read-only, so it runs under a shared lock and many workers answer it
      simultaneously.  This is safe only because of a structural property of
      the update rule, which we state and prove as \cref{thm:coverage}: the
      covered region never shrinks, so a ``covered'' verdict cannot go stale.

\item \textbf{Make the critical section small in the first place.}  Two
      mechanisms do this.  The \emph{convex-union merge}
      (\cref{sec:antichain}) collapses each discrete state's antichain to
      $\approx 1$ zone on every process network we measure, so the loop inside
      the lock is $O(1)$ rather than $O(\text{chain})$.  The
      \emph{minimal-graph inclusion test} (\cref{sec:storage}) answers the
      dominant case directly on the compressed representation, so the lock is
      not held across a cubic decompression.
\end{enumerate}

The rest of this document develops these in order: \cref{sec:background} sets
up just enough zone theory; \cref{sec:antichain} explains the sequential
insertion procedure that everything else has to preserve;
\cref{sec:parallel} is the parallel architecture; \cref{sec:storage} is the
storage layer that makes it fit in cache; \cref{sec:theory} proves the parts
that need proving; \cref{sec:uppaal} compares against \uppaal{} and \ltsmin{};
\cref{sec:eval} reports what was measured; \cref{sec:novelty} draws the line
between new and standard.
